\documentclass[11pt,a4paper]{article}

\usepackage[T1]{fontenc}
\usepackage[utf8]{inputenc}
\usepackage{amsmath,amssymb}
\usepackage[hidelinks]{hyperref}

\hypersetup{
  pdftitle={Bias, Variance, Model Complexity, and Double Descent},
  pdfauthor={Bibliographic synthesis for TP 1}
}

\newcommand{\ED}{\mathbb{E}_{\mathcal D}}
\newcommand{\VarD}{\operatorname{Var}_{\mathcal D}}
\newcommand{\EX}{\mathbb{E}_{X}}
\newcommand{\Eeps}{\mathbb{E}_{\varepsilon}}
\newcommand{\R}{\mathcal R}
\newcommand{\ind}{\mathbf{1}}

\title{\textbf{Bias, Variance, Model Complexity, and Double Descent}\\
\large A self-contained bibliographic synthesis for the practical assignment}
\author{}
\date{September 2026}

\begin{document}
\maketitle

\begin{abstract}
This document consolidates the material needed for the practical assignment on bias and variance. It integrates the classical statistical account of Hastie, Tibshirani, and Friedman; the nonparametric and neural-network perspective of Geman, Bienenstock, and Doursat; the modern double-descent account of Belkin and collaborators; and the practical model-selection viewpoint of the scikit-learn documentation. The discussion is self-contained: probability spaces and sources of randomness are stated explicitly, the squared-error decomposition is derived step by step, the special difficulties of zero-one classification loss are explained, and the theory is connected directly to polynomial regression, nearest-neighbor regression, cross-validation, heteroscedastic noise, regularization, and over-parameterized models.
\end{abstract}

\tableofcontents

\section{What this document is for}

The assignment has two related goals. The first is mathematical: understand why expected squared prediction error can be decomposed into irreducible noise, squared bias, and variance. The second is methodological: understand when this decomposition helps in model selection and where the familiar ``more complexity means more variance'' story becomes incomplete.

The central conclusions are:
\begin{enumerate}
  \item The bias-variance formula is an exact algebraic identity for squared loss under clearly stated assumptions. It is not merely a heuristic.
  \item Bias and variance are properties of a \emph{learning procedure over repeated training sets}, not properties of one fitted curve viewed in isolation.
  \item The usual additive formula does not transfer unchanged to zero-one classification loss, because thresholding probabilities into class labels is nonlinear.
  \item Model flexibility often lowers bias and raises variance in the classical regime, but raw parameter count is not a universal measure of effective complexity.
  \item Double descent extends the classical U-shaped test-risk curve beyond the interpolation threshold. It does not invalidate the squared-error identity; it invalidates an overly simple monotone story about how variance or test risk must depend on parameter count.
\end{enumerate}

\section{Statistical setup and sources of randomness}
\label{sec:setup}

\subsection{Population, training set, and learning algorithm}

Let $(X,Y)$ be a random input-output pair drawn from an unknown population distribution $P$. In regression, assume
\begin{equation}
  Y=f(X)+\varepsilon,
  \qquad
  \mathbb E[\varepsilon\mid X]=0,
  \qquad
  \operatorname{Var}(\varepsilon\mid X=x)=\sigma^2(x).
  \label{eq:data-model}
\end{equation}
The regression function is therefore
\begin{equation}
  f(x)=\mathbb E[Y\mid X=x].
\end{equation}
In the homoscedastic case, $\sigma^2(x)=\sigma^2$ is constant. In the heteroscedastic case it changes with $x$.

A training set of size $n$ is
\begin{equation}
  \mathcal D=\{(X_1,Y_1),\ldots,(X_n,Y_n)\},
\end{equation}
usually modeled as an independent sample from $P$. A learning algorithm $A$ maps the random dataset to a fitted predictor,
\begin{equation}
  \widehat f_{\mathcal D}=A(\mathcal D).
\end{equation}
Even if $A$ is deterministic, $\widehat f_{\mathcal D}$ is random because $\mathcal D$ is random. If the algorithm itself uses randomness, such as random initialization, stochastic feature selection, or bootstrap resampling, that randomness can either be included in $\mathcal D$ symbolically or represented by a separate random variable. The practical assignment mainly studies randomness caused by drawing new training sets.

It is useful to distinguish a \emph{learning procedure} from one fitted model. The procedure is the complete rule ``receive a dataset of size $n$, preprocess it, choose any hyperparameters, and fit a predictor.'' Before data are collected, its output $\widehat f_{\mathcal D}$ is random. After one observed dataset $\mathcal D_{\mathrm{obs}}$ has been supplied, the fitted function $\widehat f_{\mathcal D_{\mathrm{obs}}}$ is fixed, apart from any remaining algorithmic randomness. Bias and variance describe the procedure over possible datasets, not merely the curve obtained from the one dataset that happened to be observed.

\subsection{Loss, empirical error, and risk}

A loss function measures the cost of one prediction. Under squared loss,
\begin{equation}
  L_{2}(y,\widehat y)=(y-\widehat y)^2,
\end{equation}
whereas binary zero-one loss is
\begin{equation}
  L_{01}(y,g(x))=\ind\{y\neq g(x)\}.
\end{equation}
The latter is either zero or one for a single observation. A \emph{risk} is an expected loss under a population distribution. For example,
\begin{equation}
  \R_{01}(g)
  =\mathbb E_{X,Y}[L_{01}(Y,g(X))]
  =\Pr(Y\neq g(X)).
\end{equation}
Thus loss concerns one outcome, while risk concerns the long-run average over possible outcomes. From a finite sample one computes an empirical risk,
\begin{equation}
  \widehat{\R}_{01}(g)
  =\frac{1}{n}\sum_{i=1}^{n}L_{01}(Y_i,g(X_i)),
\end{equation}
which estimates the unknown population risk. The phrase \emph{test error} often denotes the same finite-sample average on a held-out test set; in informal writing it is sometimes also used for the population risk, so the surrounding expectation should always be checked.

Risk can be conditional or integrated. At a fixed input $x$,
\begin{equation}
  \R(g\mid x)=\mathbb E[L(Y,g(x))\mid X=x].
\end{equation}
Averaging this quantity over future inputs gives the integrated risk
\begin{equation}
  \R(g)=\mathbb E_X[\R(g\mid X)].
\end{equation}
This distinction between a lowercase fixed $x$ and a random uppercase $X$ recurs throughout the document.

\subsection{Three different expectations that must not be confused}

An expectation is shorthand for a repeated random experiment. The subscript under $\mathbb E$ states which parts of that experiment are repeated. Three random mechanisms may appear:
\begin{description}
  \item[Training-set randomness:] draw a new dataset $\mathcal D$ of the same fixed size from the same population, repeat every data-dependent modeling step, and evaluate the resulting $\widehat f_{\mathcal D}(x)$. Bias and model variance use this distribution.
  \item[Test-noise randomness:] even at a fixed $x$, a new response is $Y=f(x)+\varepsilon_{\mathrm{test}}$. A fresh $\varepsilon_{\mathrm{test}}$ creates irreducible outcome uncertainty and must be independent of the training noises.
  \item[Test inputs:] for population-wide performance, also draw $X_{\mathrm{test}}\sim P_X$. A pointwise calculation at one fixed $x$ does not include this average.
\end{description}

Consider the assignment's procedure ``draw $n=30$ observations and fit a degree-five polynomial.'' At $x=0.7$, repeated training sets might produce predictions $1.10$, $0.91$, $1.27$, and $1.02$. Their center and spread concern randomness in $\mathcal D$. A new response at the same input is $f(0.7)+\varepsilon_{\mathrm{test}}$; its additional fluctuation concerns test noise. If the application asks for average performance over future individuals, $0.7$ is no longer fixed and a new $X_{\mathrm{test}}$ is drawn as well.

The same probability model applies to real data even though the repetitions are not physically available. Before collection, another sample of $n$ individuals would have produced another dataset and another fitted predictor. After the actual dataset $\mathcal D_{\mathrm{obs}}$ has been observed and the model has been fitted, one may instead condition on that fixed dataset and study only future $(X_{\mathrm{test}},Y_{\mathrm{test}})$. Simulations make all of these repetitions observable; real applications generally provide only one realization of the training set.

The relevant random quantities can therefore be summarized as follows:
\begin{description}
  \item[Pointwise bias and variance:] fix $x$ and $f$; redraw $\mathcal D$ and refit the complete procedure.
  \item[Pointwise total test error:] fix $x$ and $f$; redraw $\mathcal D$ and an independent $\varepsilon_{\mathrm{test}}$.
  \item[Integrated risk of the procedure:] redraw $\mathcal D$, $X_{\mathrm{test}}$, and the corresponding test response noise.
  \item[Risk of one already fitted model:] condition on $\mathcal D_{\mathrm{obs}}$; average only over future test inputs and responses.
\end{description}

Hastie et al. use $\mathcal T$ for the training set in parts of Chapter 7, whereas the assignment and this document use $\mathcal D$. These are not different concepts. Geman et al. write a fitted rule as $f(x;\mathcal D)$; here it is written $\widehat f_{\mathcal D}(x)$ to distinguish the unknown truth $f$ from its estimate.

\section{The squared-error decomposition}
\label{sec:decomposition}

\subsection{Definitions at a fixed input}

Fix $X=x$. Imagine drawing many independent training sets of the same size and applying exactly the same learning procedure to each one. Even though the input $x$ is unchanged, the predictions generally differ:
\begin{equation*}
  \widehat f_{\mathcal D_1}(x),\quad
  \widehat f_{\mathcal D_2}(x),\quad\ldots
\end{equation*}
For example, at $x=0.7$ repeated degree-five fits might give $1.10$, $0.91$, $1.27$, and $1.02$. These values form a sampling distribution of the fitted prediction. Define its center by
\begin{equation}
  m(x)=\ED\left[\widehat f_{\mathcal D}(x)\right].
\end{equation}
The pointwise bias is
\begin{equation}
  \operatorname{Bias}(x)
  =m(x)-f(x),
\end{equation}
and the pointwise variance is
\begin{equation}
  \operatorname{Var}(x)
  =\VarD\left(\widehat f_{\mathcal D}(x)\right)
  =\ED\left[\left(\widehat f_{\mathcal D}(x)-m(x)\right)^2\right].
\end{equation}
Bias can be positive or negative, while squared bias and variance are nonnegative.

These definitions separate two questions. First, is the center of the fitted predictions close to the truth? This is measured by bias. Second, do the fitted predictions change substantially when the sample changes? This is measured by variance. A procedure can be stably wrong, producing high bias and low variance, or centered correctly but highly unstable, producing low bias and high variance.

Model variance must not be confused with outcome-noise variance:
\begin{equation}
  \underbrace{\VarD(\widehat f_{\mathcal D}(x))}_{\text{variation across fitted models}}
  \neq
  \underbrace{\operatorname{Var}(Y\mid X=x)}_{\text{variation of outcomes at fixed }x}.
\end{equation}
The first could in principle be reduced by a more stable learning procedure or more data. The second remains even if the regression function $f$ were known exactly.

\subsection{Complete derivation}

For a new response $Y=f(x)+\varepsilon$, independent of the training set conditional on $x$, the expected pointwise squared prediction error is
\begin{equation}
  \operatorname{Err}(x)
  =\mathbb E_{\mathcal D,\varepsilon}
   \left[\left(Y-\widehat f_{\mathcal D}(x)\right)^2\mid X=x\right].
\end{equation}
Substitute the data-generating model and add and subtract $m(x)$:
\begin{align}
  Y-\widehat f_{\mathcal D}(x)
  &=f(x)+\varepsilon-\widehat f_{\mathcal D}(x)\\
  &=\underbrace{f(x)-m(x)}_{\text{systematic error}}
    +\underbrace{m(x)-\widehat f_{\mathcal D}(x)}_{\text{training-set fluctuation}}
    +\underbrace{\varepsilon}_{\text{test noise}}.
\end{align}
Squaring gives three squared terms and three cross terms:
\begin{align}
  \left(Y-\widehat f_{\mathcal D}(x)\right)^2
  &=\left(f(x)-m(x)\right)^2
    +\left(m(x)-\widehat f_{\mathcal D}(x)\right)^2
    +\varepsilon^2 \nonumber\\
  &\quad+2\left(f(x)-m(x)\right)
             \left(m(x)-\widehat f_{\mathcal D}(x)\right) \nonumber\\
  &\quad+2\left(f(x)-m(x)\right)\varepsilon
    +2\left(m(x)-\widehat f_{\mathcal D}(x)\right)\varepsilon.
\end{align}
Every cross term has expectation zero:
\begin{align}
  \ED[m(x)-\widehat f_{\mathcal D}(x)]&=0,\\
  \mathbb E[\varepsilon\mid X=x]&=0,\\
  \mathbb E[(m(x)-\widehat f_{\mathcal D}(x))\varepsilon\mid X=x]&=0,
\end{align}
where the last equality uses conditional independence between a new test noise and the training set. Therefore
\begin{equation}
  \boxed{
  \operatorname{Err}(x)
  =\underbrace{\left(f(x)-m(x)\right)^2}_{\operatorname{Bias}^2(x)}
   +\underbrace{\VarD\left(\widehat f_{\mathcal D}(x)\right)}_{\operatorname{Variance}(x)}
   +\underbrace{\sigma^2(x)}_{\text{irreducible noise}}.}
  \label{eq:bv-pointwise}
\end{equation}

The noise is called irreducible because even the oracle predictor $\widehat f(x)=f(x)$ has expected error
\begin{equation}
  \mathbb E[(Y-f(x))^2\mid X=x]=\sigma^2(x).
\end{equation}
It is irreducible relative to predicting the random response $Y$ with a deterministic point prediction. It may still be scientifically useful to model $\sigma^2(x)$, for example to construct prediction intervals.

\subsection{Integrated or average risk}

If test inputs are also random, average 
equation~\eqref{eq:bv-pointwise} over $P_X$:
\begin{equation}
  \mathbb E_X[\operatorname{Err}(X)]
  =\mathbb E_X[\operatorname{Bias}^2(X)]
   +\mathbb E_X[\operatorname{Variance}(X)]
   +\mathbb E_X[\sigma^2(X)].
  \label{eq:integrated}
\end{equation}
The distribution used for the $X$ average matters. A uniform test grid approximates a uniform integral over $[0,1]$; it does not approximate an average under a Beta$(2,5)$ test distribution unless weights are added. This distinction becomes important in Exercise 6.

\subsection{Model bias and estimation bias}

Hastie et al. further separate two reasons for systematic error. Suppose the chosen model class is $\mathcal H$. Define its best population approximation
\begin{equation}
  f_{\mathcal H}^{*}
  =\arg\min_{h\in\mathcal H}\mathbb E_X[(f(X)-h(X))^2].
\end{equation}
Then one can conceptually distinguish
\begin{itemize}
  \item \emph{model or approximation bias}: $f-f_{\mathcal H}^{*}$, caused by a model class that cannot represent the truth;
  \item \emph{estimation bias}: $f_{\mathcal H}^{*}-\ED[\widehat f_{\mathcal D}]$, caused by restrictions or penalties in the fitting procedure.
\end{itemize}
Under suitable orthogonality conditions, as in population least-squares projection, their integrated squared contributions add. In general, expanding a square also produces a cross term, so the conditions behind any claimed additive sub-decomposition should be checked.

For example, fitting a straight line to a strongly oscillatory truth creates approximation bias because no member of the linear class can reproduce the oscillations. Fitting a sufficiently rich polynomial with a strong ridge penalty can create estimation bias even if the unpenalized class contains a good approximation: the fitting rule deliberately shrinks the solution away from the best unpenalized coefficients. The ordinary bias $m(x)-f(x)$ combines both effects unless an additional decomposition and its assumptions are stated.

\section{How the decomposition is estimated in the practical assignment}
\label{sec:monte-carlo}

The practical assignment makes the otherwise hypothetical repetitions observable because the data-generating function is known. Conceptually, the computation follows seven steps:
\begin{enumerate}
  \item draw $M$ independent training sets of the same size;
  \item refit the chosen learning procedure on every dataset;
  \item evaluate every fitted model at the same $J$ test inputs;
  \item at each input, compute the center and spread of the $M$ predictions;
  \item compare the center with the known $f(x)$ to estimate squared bias;
  \item generate fresh noisy test responses, independently of all training noises, to estimate total test error;
  \item average pointwise quantities over the requested test-input distribution.
\end{enumerate}
This sequence should be kept conceptually separate even if an implementation combines several operations in arrays.

Suppose $M$ independent training sets produce predictions
\begin{equation}
  \widehat f^{(1)}(x_j),\ldots,\widehat f^{(M)}(x_j)
\end{equation}
at test-grid points $x_1,\ldots,x_J$. The Monte Carlo mean prediction is
\begin{equation}
  \overline f(x_j)=\frac{1}{M}\sum_{m=1}^{M}\widehat f^{(m)}(x_j).
\end{equation}
Natural empirical estimates are
\begin{align}
  \widehat{\operatorname{Bias}^2}(x_j)
  &=\left(\overline f(x_j)-f(x_j)\right)^2,\\
  \widehat{\operatorname{Var}}(x_j)
  &=\frac{1}{M}\sum_{m=1}^{M}
    \left(\widehat f^{(m)}(x_j)-\overline f(x_j)\right)^2.
  \label{eq:mcvar}
\end{align}
Using $1/(M-1)$ would make the usual sample variance unbiased for the variance of a scalar random variable. Using $1/M$ instead matches the direct Monte Carlo average appearing in the theoretical definition. The difference is negligible for $M=200$, but the convention should be reported consistently.

Averaging over a uniform grid yields
\begin{align}
  \widehat{B^2}&=\frac{1}{J}\sum_{j=1}^{J}
      \widehat{\operatorname{Bias}^2}(x_j),\\
  \widehat V&=\frac{1}{J}\sum_{j=1}^{J}
      \widehat{\operatorname{Var}}(x_j).
\end{align}
The total test error is the expected squared error on a new response, averaged over the random training set, the test input, and a fresh test noise:
\begin{align*}
  E_{\mathrm{test}}
  &=\mathbb E_{\mathcal D,X,\varepsilon}
    \left[\left(Y-\widehat f_{\mathcal D}(X)\right)^2\right]\\
  &=\mathbb E_{\mathcal D,X,\varepsilon}
    \left[\left(f(X)+\varepsilon-
    \widehat f_{\mathcal D}(X)\right)^2\right].
\end{align*}
This expression contains three expectations. Each is replaced by a finite Monte Carlo average in the experiment.

First, the expectation over possible training sets is approximated by fitting $M$ models on $M$ independently drawn datasets:
\begin{align*}
  E_{\mathrm{test}}
  &\approx \frac{1}{M}\sum_{m=1}^{M}
  \mathbb E_{X,\varepsilon}
  \left[\left(f(X)+\varepsilon-
  \widehat f^{(m)}(X)\right)^2\right].
\end{align*}
Second, the expectation over a uniformly distributed test input is approximated by the arithmetic mean over the $J$ points of the uniform grid:
\begin{align*}
  E_{\mathrm{test}}
  &\approx \frac{1}{MJ}\sum_{m=1}^{M}\sum_{j=1}^{J}
  \mathbb E_{\varepsilon}
  \left[\left(f(x_j)+\varepsilon-
  \widehat f^{(m)}(x_j)\right)^2\right].
\end{align*}
Finally, for every model-grid pair $(m,j)$, draw an independent fresh test noise
\[
  \varepsilon_j^{(m)}\sim\mathcal N(0,\sigma^2)
\]
and define the corresponding new test response as
\[
  Y_j^{(m)}=f(x_j)+\varepsilon_j^{(m)}.
\]
Replacing the remaining noise expectation by this Monte Carlo draw gives
\begin{equation}
  \widehat E_{\mathrm{test}}
  =\frac{1}{MJ}\sum_{m=1}^{M}\sum_{j=1}^{J}
  \left(f(x_j)+\varepsilon_j^{(m)}-\widehat f^{(m)}(x_j)\right)^2.
  \label{eq:mc-test-error}
\end{equation}
Thus, equation~\eqref{eq:mc-test-error} is simply the mean of the $M\times J$ squared errors between fresh noisy test responses and their corresponding predictions. The noise must be fresh and independent of the training noise. If it were omitted, the quantity would estimate only squared bias plus variance and would not include irreducible noise.

The empirical identity will not hold exactly because all quantities are Monte Carlo estimates. A useful signed relative discrepancy is
\begin{equation}
  \Delta_{\mathrm{rel}}
  =\frac{\widehat E_{\mathrm{test}}-
    (\widehat{B^2}+\widehat V+\sigma^2)}
    {\widehat E_{\mathrm{test}}},
\end{equation}
and an unsigned reporting measure is $|\Delta_{\mathrm{rel}}|$. The discrepancy should generally shrink when $M$ and the number of independent test-noise draws increase.

\section{Complexity, underfitting, and overfitting}
\label{sec:complexity}

\subsection{What capacity means}

Model \emph{capacity}, \emph{flexibility}, or \emph{expressiveness} refers to the range of input-output relations that a learning procedure can represent or select. A low-capacity class contains relatively restricted functions; a high-capacity class can reproduce a wider variety of patterns, including irregular ones. Capacity is therefore not synonymous with quality: the extra possibilities may represent genuine structure or sample-specific noise.

The operational capacity parameter depends on the model:
\begin{itemize}
  \item for polynomial regression, increasing the degree increases capacity;
  \item for $k$-nearest neighbors, decreasing $k$ increases capacity because predictions become more local;
  \item for ridge regression with a fixed feature set, decreasing the penalty $\alpha$ increases effective capacity;
  \item for a decision tree, increasing permitted depth or the number of leaves increases capacity.
\end{itemize}
Raw parameter count is only one possible proxy. The size of the available function class, the complexity of the particular fitted solution, and the effective capacity after regularization are different objects. This distinction becomes essential beyond the interpolation threshold.

\subsection{The classical picture}

In the classical regime, increasing flexibility often has opposing effects:
\begin{itemize}
  \item The model can represent more shapes, so approximation bias tends to decrease.
  \item The fitted predictor becomes more sensitive to the particular training sample, so variance tends to increase.
\end{itemize}
Their sum, plus irreducible noise, often produces a U-shaped test-error curve. Training error, by contrast, normally decreases as flexibility increases and is therefore not a reliable model-selection criterion.

This is a tendency, not a theorem that every bias curve must decrease monotonically or every variance curve must increase monotonically for every algorithm and dataset. The exact quantities depend on the estimator, training procedure, design distribution, sample size, and test location.

\subsection{Polynomial regression}

For polynomial regression, degree is a simple flexibility parameter. A low degree may be unable to approximate an oscillatory $f$, producing high bias and relatively stable fitted curves. A high degree can interpolate or nearly interpolate a small noisy sample; small changes in the sample may then cause large changes between observed points, especially near boundaries. That is high variance.

To make the phrase ``fit the noise'' precise, write the observed responses as
\begin{equation}
  Y_i=f(X_i)+\varepsilon_i.
\end{equation}
A sufficiently flexible interpolating polynomial satisfies
\begin{equation}
  \widehat f_{\mathcal D}(X_i)=Y_i=f(X_i)+\varepsilon_i
\end{equation}
at every training input. If one observation is unusually high only because its realized $\varepsilon_i$ is positive, the fitted curve may create a local rise to reach it. A new dataset contains different noises and can therefore produce a very different high-degree curve. The resulting sensitivity across datasets is precisely model variance.

A lower-degree polynomial cannot follow every displacement. It may leave nonzero training residuals while capturing the common trend $f$. This can be beneficial because the particular training noises do not repeat in future data. Nevertheless, interpolation alone is not a proof of poor generalization: in a highly over-parameterized regime there may be many interpolating solutions, and the algorithm used to select among them matters. Section~\ref{sec:double-descent} develops this qualification.

Monomial bases at high degree are often numerically ill conditioned. An extreme curve may therefore reflect both statistical variance and numerical instability. Rescaling inputs, using an orthogonal polynomial basis, or applying ridge regularization can improve numerical behavior. Such changes do not alter the conceptual decomposition, but they can alter the actual learning algorithm and therefore its measured bias and variance.

\subsection{Nearest-neighbor regression}

For $k$-nearest-neighbor regression,
\begin{equation}
  \widehat f_k(x)=\frac{1}{k}\sum_{i\in N_k(x)}Y_i.
\end{equation}
Here complexity runs in the opposite direction from $k$:
\begin{itemize}
  \item Small $k$: local, flexible fit; low smoothing; typically lower bias and higher variance.
  \item Large $k$: broad averaging; strong smoothing; typically higher bias and lower variance.
\end{itemize}
If the selected neighbors and their input locations are treated as fixed and their noises are independent with variance $\sigma^2$, then
\begin{equation}
  \operatorname{Var}(\widehat f_k(x)\mid X_1,\ldots,X_n)
  =\frac{\sigma^2}{k}.
\end{equation}
This follows because the average contains $k$ independent noise terms:
\begin{equation*}
  \operatorname{Var}\left(\frac{1}{k}\sum_{i\in N_k(x)}\varepsilon_i\right)
  =\frac{1}{k^2}\sum_{i\in N_k(x)}\operatorname{Var}(\varepsilon_i)
  =\frac{k\sigma^2}{k^2}.
\end{equation*}
With $k=1$, the prediction inherits all the selected neighbor's noise. With $k=20$, independent noise is strongly averaged, but the neighborhood may extend far enough that values of $f$ different from $f(x)$ are mixed together. This is the local form of the bias-variance trade-off.

The conditional bias is approximately the difference between $f(x)$ and the average of $f$ over the neighbor locations. With random training inputs, additional variability comes from the random neighbor locations themselves.

\subsection{Why more data usually reduces variance}

Increasing $n$ supplies more information about the same population. Parameter estimates typically become more stable, and local methods can use enough nearby points without smoothing across a large region. Therefore a more flexible model may become optimal as $n$ grows.

This statement is not automatic under distribution shift or severe misspecification. More observations from the wrong region of input space do not necessarily reduce error where data remain scarce.

\section{The nonparametric perspective of Geman, Bienenstock, and Doursat}

Geman et al. interpret feedforward neural networks as nonparametric regression estimators. Their central dilemma is about learning a flexible function from finite data:
\begin{itemize}
  \item A strongly restricted model encodes assumptions and can be stable, but wrong assumptions produce bias.
  \item A nearly model-free estimator can approximate many functions, but estimating such freedom from a finite sample produces variance.
\end{itemize}

They emphasize that this is fundamentally a sample-size problem. Faster hardware can optimize a flexible model more quickly, but computation alone does not create information absent from the training sample.

For example, suppose only a few observations lie in a neighborhood of $x$. A highly local estimator is being asked to infer detailed structure from very little information, so its answer can change substantially when those few observations change. Increasing computation does not reveal the unobserved local responses. One must instead collect more relevant data or impose structure that lets information be shared across locations.

Their term \emph{smoothing parameter} covers quantities such as the number of neighbors, kernel bandwidth, and, in their neural-network discussion, the number of hidden units. Smoothing introduces regularity and thus some bias to control variance. They discuss cross-validation as a data-driven method for selecting smoothing and regularization as a way of encoding a preference for particular functions.

Smoothing deliberately lets nearby or structurally related inputs influence one another. Increasing $k$ in $k$-NN averages more responses; increasing a kernel bandwidth averages over a wider region; increasing a regularization penalty suppresses rapidly changing or weakly supported directions. All three reduce sensitivity to individual observations, but can also erase genuine fine-scale structure. The useful amount of smoothing depends on both sample size and the regularity of the unknown function.

An important historical qualification is needed. The paper predates modern large-scale deep learning. Its claim that difficult learning problems require strong prior representations remains influential, but modern systems may acquire useful representations from large datasets, pretraining, architecture, optimization, and augmentation. These mechanisms are forms of inductive bias; they do not constitute learning without assumptions.

\section{Classification and zero-one loss}
\label{sec:classification}

\subsection{Probability estimation versus class prediction}

For binary classification, let $Y\in\{0,1\}$ and
\begin{equation}
  \eta(x)=\Pr(Y=1\mid X=x).
\end{equation}
A learning algorithm may first estimate $\eta(x)$ and then classify by
\begin{equation}
  \widehat g_{\mathcal D}(x)
  =\ind\{\widehat\eta_{\mathcal D}(x)\geq 1/2\}.
\end{equation}
The squared error of the probability estimator $\widehat\eta$ has the standard additive bias-variance decomposition. However, the classification error uses zero-one loss,
\begin{equation}
  L_{01}(Y,\widehat g)=\ind\{Y\neq\widehat g\},
\end{equation}
which depends only on which side of the threshold the prediction lies.

The notation distinguishes an individual loss from its expectation. For one realized label $y$, $L_{01}(y,g(x))$ is either zero or one. The conditional risk
\begin{equation}
  \R_{01}(g\mid x)
  =\mathbb E[L_{01}(Y,g(x))\mid X=x]
\end{equation}
is the probability of error among all possible observations having the same input $x$. The integrated risk
\begin{equation}
  \R_{01}(g)=\mathbb E_X[\R_{01}(g\mid X)]
\end{equation}
is the population-wide probability of misclassification. A finite test-set error estimates this last quantity but is not identical to the unknown population expectation.

\subsection{Conditional classification risk and the Bayes classifier}

For any deterministic class prediction $g(x)\in\{0,1\}$,
\begin{align}
  \R_{01}(g\mid x)
  &=\Pr(Y\neq g(x)\mid X=x)\\
  &=\eta(x)\ind\{g(x)=0\}
    +(1-\eta(x))\ind\{g(x)=1\}.
\end{align}
The second equality is simply a compact representation of two cases. If $g(x)=0$, an error occurs exactly when $Y=1$, whose conditional probability is $\eta(x)$. If $g(x)=1$, an error occurs when $Y=0$, whose probability is $1-\eta(x)$. Thus
\begin{equation*}
  \R_{01}(g\mid x)
  =\begin{cases}
    \eta(x), & g(x)=0,\\
    1-\eta(x), & g(x)=1.
  \end{cases}
\end{equation*}
This is minimized by the Bayes classifier
\begin{equation}
  g^{*}(x)=\ind\{\eta(x)\geq 1/2\},
\end{equation}
whose irreducible conditional classification error is
\begin{equation}
  \R_{01}^{*}(x)=\min\{\eta(x),1-\eta(x)\}.
\end{equation}
Unlike Gaussian regression noise, this quantity is not generally denoted $\sigma^2$.

For example, if $\eta(x)=0.8$, predicting class 1 has conditional error $0.2$, whereas predicting class 0 has error $0.8$. Bayes therefore selects class 1, but it still makes an error in the $20\%$ of outcomes for which $Y=0$. This remaining $0.2$ is irreducible for deterministic classification using only $x$. If $\eta(x)=0.5$, neither class is preferable and the minimum conditional error is $0.5$.

The excess conditional zero-one risk has a particularly informative form:
\begin{equation}
  \R_{01}(g\mid x)-\R_{01}^{*}(x)
  =|2\eta(x)-1|\,\ind\{g(x)\neq g^{*}(x)\}.
  \label{eq:excess01}
\end{equation}
To verify it, consider separately $\eta(x)\geq 1/2$ and $\eta(x)<1/2$. If $g$ agrees with the Bayes decision, the difference is zero. If it chooses the opposite class, the two conditional risks differ by $|2\eta(x)-1|$.

The verification can be written explicitly. When $\eta(x)\geq 1/2$, Bayes predicts 1 and has risk $1-\eta(x)$. If $g$ instead predicts 0, then
\begin{equation*}
  \R_{01}(g\mid x)-\R_{01}^{*}(x)
  =\eta(x)-(1-\eta(x))
  =2\eta(x)-1.
\end{equation*}
When $\eta(x)<1/2$, Bayes predicts 0 and has risk $\eta(x)$. If $g$ instead predicts 1, then
\begin{equation*}
  \R_{01}(g\mid x)-\R_{01}^{*}(x)
  =(1-\eta(x))-\eta(x)
  =1-2\eta(x).
\end{equation*}
Both cases equal $|2\eta(x)-1|$. The indicator in equation~\eqref{eq:excess01} makes this cost zero whenever $g$ agrees with Bayes.

The factor
\begin{equation}
  |2\eta(x)-1|=2|\eta(x)-1/2|
\end{equation}
is a decision margin. If $\eta(x)=0.51$, choosing the wrong side of the threshold increases conditional risk by only $0.02$. If $\eta(x)=0.90$, the same disagreement increases risk by $0.80$. Errors near an intrinsically ambiguous boundary are therefore less costly than errors where one class is overwhelmingly more likely.

\subsection{Why the regression identity does not carry over}

The derivation in Section~\ref{sec:decomposition} works because a square can be expanded and centered, making cross terms vanish. The indicator loss has no analogous quadratic expansion. Two consequences follow:
\begin{enumerate}
  \item A probability estimate can have sizable squared bias without changing the class decision. If $\eta(x)=0.9$ and the mean estimate is $0.6$, the probability estimate is biased, but both values imply class 1.
  \item Variability matters mainly when it causes predictions to cross the decision boundary. The interaction between the mean prediction, its dispersion, and the distance from $1/2$ is essential.
\end{enumerate}
Thus the standard statement
\begin{equation*}
  \text{zero-one error}=\text{noise}+\text{squared bias}+\text{variance}
\end{equation*}
is false if ``bias'' and ``variance'' retain their squared-loss definitions.

Generalized decompositions for arbitrary losses do exist. Domingos defines a loss-dependent central prediction and corresponding loss-based bias and variance. Under squared loss, the central prediction is the mean and the standard formula is recovered; under zero-one loss, the central prediction is the mode. These generalized quantities are useful, but they should not be silently identified with the ordinary squared bias and prediction variance. This resolves the apparent disagreement: Hastie et al. reject the unchanged squared-loss sum for zero-one error, while Domingos changes the definitions to construct a loss-specific decomposition.

\section{Regularization and practical levers}
\label{sec:levers}

\subsection{Levers that usually reduce variance}

\begin{itemize}
  \item Increase the training-set size, especially in regions where predictions are unstable.
  \item Reduce model flexibility: lower polynomial degree, increase $k$ in $k$-NN, prune a tree, or reduce the number of active features.
  \item Strengthen explicit regularization, such as increasing the ridge penalty.
  \item Average unstable predictors through bagging or other ensembles, provided their errors are not perfectly correlated.
  \item Use early stopping or data augmentation when these mechanisms impose a suitable effective regularity.
\end{itemize}

These interventions may increase bias. The goal is not to minimize variance alone, but to minimize total generalization error.

\subsection{Levers that usually reduce bias}

\begin{itemize}
  \item Use a richer hypothesis class: increase polynomial degree, add relevant transformations or interactions, decrease $k$ in $k$-NN, or use a more expressive architecture.
  \item Weaken regularization when it prevents the model from capturing genuine structure.
  \item Improve features or representation so that the target function becomes simpler in the chosen coordinates.
  \item Optimize sufficiently well. Insufficient training can cause systematic error that resembles statistical bias.
\end{itemize}

Reducing bias may increase variance. The effect of any lever must therefore be checked on independent data or by a correctly nested validation procedure.

\subsection{Ridge regression as an explicit trade-off}

For a design matrix $X$ and response vector $y$, ridge regression solves
\begin{equation}
  \widehat\beta_{\alpha}
  =\arg\min_{\beta}
   \left\{\|y-X\beta\|_2^2+\alpha\|\beta\|_2^2\right\}
  =(X^{\mathsf T}X+\alpha I)^{-1}X^{\mathsf T}y.
\end{equation}
At a test vector $x_0$, the prediction is linear in $y$:
\begin{equation}
  \widehat f_{\alpha}(x_0)
  =x_0^{\mathsf T}(X^{\mathsf T}X+\alpha I)^{-1}X^{\mathsf T}y.
\end{equation}
Increasing $\alpha$ shrinks directions associated with small eigenvalues of $X^{\mathsf T}X$. This typically lowers variance, especially for ill-conditioned polynomial features, while introducing estimation bias. The optimal $\alpha$ balances the two.

To interpret the reference to eigenvalues, imagine rotating the coefficient coordinates into directions in which $X^{\mathsf T}X$ is diagonal. If one direction has eigenvalue $\lambda_j$, ridge multiplies the corresponding unregularized component by a factor of the form
\begin{equation}
  \frac{\lambda_j}{\lambda_j+\alpha}.
\end{equation}
A small $\lambda_j$ describes a direction that the observed data constrain weakly. Its fitted coefficient can change greatly in response to small perturbations, so ridge shrinks it strongly. Well-supported directions with large $\lambda_j$ are changed less. This explains how ridge stabilizes ill-conditioned high-degree polynomial fits without literally removing the high-degree columns.

\subsection{Simplifying a flexible fitted model}

There are several meanings of ``simplify the model,'' and they should not be conflated. In the assignment, the most direct method is to compare complete polynomial models of degrees $0$ through $15$ by validation and select one degree. If degree five is selected, the final model retains the hierarchy
\begin{equation}
  1,x,x^2,\ldots,x^5.
\end{equation}
Keeping a high-order term while arbitrarily deleting lower-order terms is usually avoided unless subject-matter structure justifies it.

A second strategy is to retain a rich feature set and reduce its effective complexity through regularization. Ridge shrinks coefficients continuously but generally does not set them exactly to zero. An $\ell_1$ penalty, as in the Lasso, can set some coefficients to zero and therefore perform a form of feature selection.

Simply deleting the numerically smallest monomial coefficients is unreliable. The magnitude of $\beta_j$ depends on the scaling of $x^j$, and polynomial columns can be strongly correlated. A term such as $0.001x^{15}$ may be influential over a sufficiently large input scale, while correlated terms can distribute one effect across several coefficients. Any data-driven simplification must therefore be performed inside the validation procedure and the reduced model should be refitted and evaluated independently.

\section{Double descent and the interpolation regime}
\label{sec:double-descent}

\subsection{The extended risk curve}

Belkin et al. organize classical and modern behavior on one capacity axis:
\begin{enumerate}
  \item \textbf{Under-parameterized classical regime:} increasing capacity initially reduces test risk by lowering underfitting.
  \item \textbf{Classical overfitting region:} near the best classical capacity, further flexibility increases test risk.
  \item \textbf{Interpolation threshold:} the smallest capacity at which training data can be fitted exactly. Test risk may peak here because fitting every noisy constraint can require an unstable solution.
  \item \textbf{Over-parameterized interpolation regime:} capacity increases further while training error remains zero, yet test risk can decrease again. This is the second descent.
\end{enumerate}

The term ``double descent'' refers to the resulting test-risk curve: a first descent, a rise near interpolation, and a second descent beyond interpolation.

Consider a concrete but illustrative experiment with $n=100$ noisy training observations. A two-parameter model may have training risk $0.40$ and test risk $0.42$ because it cannot represent the signal. With twenty effective parameters, these values might fall to $0.12$ and $0.15$. Near eighty parameters, training risk might fall to $0.02$ while test risk rises to $0.25$ as the model begins to track sample-specific noise. At the smallest capacity capable of exact fitting, training risk reaches zero and test risk may peak, for example at $0.60$. Far beyond that threshold, training risk remains zero while test risk may fall again, perhaps to $0.22$ and then $0.13$. These numbers are illustrative rather than universal; their purpose is to identify the four regimes on the capacity axis.

The interpolation threshold is defined by the training constraints, not by a particular parameter count in every model. For squared-loss regression, an interpolating solution satisfies
\begin{equation}
  \widehat f(X_i)=Y_i
  \qquad\text{for every }i=1,\ldots,n,
\end{equation}
and therefore has zero training loss. Since $Y_i=f(X_i)+\varepsilon_i$, interpolation reproduces both signal and the realized training noise at the observed inputs.

\subsection{Why more parameters can lead to a simpler fitted function}

When many interpolating solutions exist, the learning algorithm must select one. In the random-feature models studied by Belkin et al., the selected solution is the interpolant with minimum coefficient norm. Increasing the number of available features can make it possible to satisfy the same interpolation constraints with a smaller norm and a smoother function. Thus:
\begin{equation*}
  \text{larger function class}
  \not\Rightarrow
  \text{more complex selected predictor}.
\end{equation*}
The class size, the effective complexity of the selected solution, and the inductive bias of the algorithm are distinct notions.

The existence of many interpolants can be seen without neural networks. Suppose a polynomial $p(x)$ passes through observations at $x=0,1,2$. Then, for any constant $c$,
\begin{equation}
  p_c(x)=p(x)+c\,x(x-1)(x-2)
\end{equation}
passes through exactly the same observations, because the added term is zero at all three training inputs. The functions can nevertheless behave very differently between and beyond those inputs. Equal training error therefore does not imply equal test behavior.

In over-parameterized linear regression the same idea appears as
\begin{equation}
  X\beta=y.
\end{equation}
When there are more parameters than independent constraints, this equation can have infinitely many solutions. An algorithm may select the minimum-norm interpolant
\begin{equation}
  \widehat\beta
  =\arg\min_{\beta:X\beta=y}\|\beta\|_2.
\end{equation}
Near the interpolation threshold, there may be just enough freedom to satisfy every noisy constraint, requiring a highly unstable solution. With much more freedom, the same constraints may be shared among many directions and a lower-norm solution may be available. This is the intuition behind the possible second descent; it is not a guarantee for arbitrary polynomial bases or optimization procedures.

For neural networks, optimization can impose implicit regularization: among many zero-training-error solutions, gradient-based training may favor particular norms, margins, or other structural properties. Early stopping adds an explicit dependence on the training trajectory. Architecture, initialization, augmentation, and preprocessing also change the effective hypothesis space.

\subsection{What double descent does and does not contradict}

Double descent does \emph{not} contradict the algebraic identity in equation~\eqref{eq:bv-pointwise}. For squared loss, the identity remains true for every degree or parameter count. What fails is the extra heuristic that variance must grow monotonically with raw capacity or that every interpolating model must generalize poorly.

It is also unsafe to claim that larger models are always better. The second descent is an empirical pattern observed for important model-algorithm-data combinations, not a universal guarantee. Noise level, sample size, regularization, optimization, and the chosen capacity axis all matter. The interpolation peak can be narrow and may be hidden by coarse parameter grids or early stopping.

\subsection{Relation to the polynomial experiment}

The assignment's degrees $0$ through $15$ primarily illustrate the classical trade-off. A high-degree polynomial may approach interpolation and become unstable, but the experiment is not automatically a demonstration of a full second descent. Claiming double descent requires exploring a capacity range beyond interpolation and observing test risk fall again while training risk remains essentially zero.

\section{Cross-validation when the true function is unknown}
\label{sec:cv}

In simulations, $f$ is known and new datasets can be generated, so true pointwise bias and variance can be approximated directly. In a real application, only one training dataset is observed and $f$ is unknown. Bias and variance remain well-defined properties of the procedure under hypothetical repeated sampling, but they are latent: the available data do not reveal them separately without additional assumptions. Model selection therefore targets total predictive risk rather than trying to minimize an unobservable component in isolation.

\subsection{A concrete five-fold example}

Suppose $n=100$ observations are available. Polynomial degrees $1$ through $15$ are candidates. Five-fold cross-validation divides the observations into five groups of approximately twenty. For a fixed degree, the procedure is:
\begin{enumerate}
  \item fit the model on folds 2--5 and evaluate it on fold 1;
  \item fit a new model on folds 1, 3, 4, and 5 and evaluate it on fold 2;
  \item repeat until every fold has served once as validation data;
  \item average the five validation losses.
\end{enumerate}
The entire sequence is repeated for every candidate degree. If degrees 3 through 7 produce estimated risks $0.124$, $0.117$, $0.115$, $0.116$, and $0.119$, ordinary cross-validation selects degree 5 because $0.115$ is smallest.

This calculation estimates total generalization error. It does not reveal that degree 5 has the smallest bias or the smallest variance separately. Under squared loss, cross-validation approximately compares the complete sum
\begin{equation*}
  \text{squared bias}+\text{model variance}+\text{irreducible noise},
\end{equation*}
with finite-sample uncertainty added by the validation procedure itself.

\subsection{Why the selected degree can be unstable}

Compare two hypothetical validation curves near their minima. In the first case, degrees 4, 5, and 6 have errors
\begin{equation*}
  0.30,\qquad 0.10,\qquad 0.28.
\end{equation*}
Degree 5 is separated clearly from its competitors, so small perturbations are unlikely to change the selection. In the second case the errors are
\begin{equation*}
  0.117,\qquad 0.115,\qquad 0.116.
\end{equation*}
The curve is nearly flat. A new dataset or another assignment of observations to folds can easily change which value is smallest, even though the predictive performance of all three degrees is practically the same.

Consequently, a sequence of repeated experiments might select degrees
\begin{equation*}
  5,\ 4,\ 6,\ 5,\ 7,\ 4,\ 6,\ldots
\end{equation*}
without producing comparably large changes in test error. Hyperparameter instability and predictive instability are different claims. Typical causes of selection variability include a flat validation curve, small $n$, large noise, genuinely similar candidate risks, randomness in fold construction, and occasional numerical instability of high-degree fits. Exercise 5 should therefore report both the frequency of selected degrees and the validation-error curves or achieved errors when possible.

\subsection{Every learned step belongs inside the folds}

A validation fold must imitate data that were unavailable during training. Suppose one wants to select the five features most correlated with the response. The following procedure is invalid:
\begin{enumerate}
  \item use all observations, including future validation folds, to choose five features;
  \item cross-validate models using only those preselected features.
\end{enumerate}
The validation responses have already influenced which features were retained. This leakage makes the estimated risk optimistic even though the final regression coefficients were not fitted on the held-out observations.

The correct procedure repeats feature selection separately inside every training fold, fits the model using only that fold's training portion, and then applies the resulting transformation and model to the untouched validation fold. The same rule applies to supervised feature selection, normalization parameters, missing-value imputation, dimensionality reduction, and hyperparameter tuning: if a quantity is learned from data, it must be learned without access to the fold on which it will be evaluated.

Finally, cross-validation used to choose a degree is itself part of model selection. A separate test set, if available, must remain untouched until the degree and every other modeling choice have been fixed. Cross-validation answers ``which procedure should be selected?''; the final test set estimates how the selected procedure performs on genuinely unseen data.

\section{Nonuniform inputs and heteroscedastic noise}
\label{sec:hetero}

This exercise changes two different parts of the data-generating process. First, the input distribution determines where observations are available. Second, the conditional noise distribution determines how variable the response is at each location. Their effects should be analyzed separately before being combined.

\subsection{Nonuniform inputs and boundary instability}

Suppose training inputs follow Beta$(2,5)$ rather than a uniform distribution. This density is concentrated near the left part of $[0,1]$ and supplies few observations near $x=1$. A typical sample may contain many inputs around $0.1$--$0.4$ and only one or two above $0.8$.

At an interior location, nearby observations can lie on both sides of the prediction point and constrain the level, slope, and curvature of a fitted polynomial. Near the right boundary, there are no observations with $x>1$, and the Beta design may provide few observations even immediately to the left. A flexible polynomial can then rise or fall sharply near $x=1$ without paying a large training-error penalty. If a rare rightmost observation moves because of noise, the fitted boundary behavior may change dramatically.

Across repeated datasets, predictions near a well-sampled point might be $0.48$, $0.51$, and $0.49$, whereas predictions near $x=0.95$ might be $1.2$, $-0.3$, and $2.0$. The latter have much greater pointwise model variance. This is the meaning of the statement that flexible estimators are less constrained near the right boundary. It is a property of the sampling design and fitted procedure, not yet a statement about the response-noise variance.

\subsection{Heteroscedastic outcome noise}

Now let
\begin{equation}
  \varepsilon(x)\sim\mathcal N\left(0,
  \left[\sigma(0.5+2x)\right]^2\right).
\end{equation}
The irreducible pointwise noise is no longer $\sigma^2$ but
\begin{equation}
  \boxed{\sigma^2(x)=\sigma^2(0.5+2x)^2.}
\end{equation}
This quantity increases strongly across the interval:
\begin{align*}
  \sigma^2(0)&=0.25\sigma^2,\\
  \sigma^2(0.5)&=2.25\sigma^2,\\
  \sigma^2(1)&=6.25\sigma^2.
\end{align*}
Thus observations close to $x=1$ are intrinsically noisier even if the regression function were known. The right side can now be difficult for two separate reasons: sparse inputs can increase model variance, while heteroscedasticity increases irreducible outcome noise. The two quantities are
\begin{equation*}
  \VarD(\widehat f_{\mathcal D}(x))
  \quad\text{and}\quad
  \operatorname{Var}(Y\mid X=x)=\sigma^2(x),
\end{equation*}
and should not be merged conceptually.

The correct verification is therefore
\begin{equation}
  \operatorname{Err}(x)
  \approx \operatorname{Bias}^2(x)
  +\operatorname{Variance}(x)
  +\sigma^2(0.5+2x)^2.
\end{equation}
\subsection{Why the test-input distribution changes the average}

Pointwise error describes one location. An integrated error requires weights over locations, and those weights encode the target population. If quantities are averaged over a uniform test grid, the corresponding average noise term is
\begin{align}
  \int_0^1 \sigma^2(0.5+2x)^2\,dx
  &=\sigma^2\int_0^1(0.25+2x+4x^2)\,dx\\
  &=\sigma^2\left(0.25+1+\frac{4}{3}\right)\\
  &=\frac{31}{12}\sigma^2.
  \label{eq:uniform-hetero-noise}
\end{align}
If test inputs instead follow Beta$(2,5)$, the correct average is
\[
  \mathbb E_{X\sim\mathrm{Beta}(2,5)}
  \left[\sigma^2(0.5+2X)^2\right],
\]
which is different. Using $\mathbb E[X]=2/7$ and $\mathbb E[X^2]=2\cdot3/(7\cdot8)=3/28$ gives
\begin{align}
  \mathbb E[\sigma^2(X)]
  &=\sigma^2\left(0.25+2\mathbb E[X]+4\mathbb E[X^2]\right)\\
  &=\sigma^2\left(\frac14+\frac47+\frac37\right)
  =\frac54\sigma^2.
\end{align}
The uniform value is approximately $2.58\sigma^2$, whereas the Beta-weighted value is $1.25\sigma^2$. The uniform average gives every part of $[0,1]$ equal weight, including the highly noisy right side. The Beta$(2,5)$ average assigns little probability to that region, so its contribution to population-wide risk is smaller. Neither average is intrinsically correct: the appropriate one depends on whether the scientific target treats locations uniformly or predicts future observations drawn from the Beta population. This is why the test distribution must always be stated.

\section{Model variance, ensembles, and predictive uncertainty}

The variance term in equation~\eqref{eq:bv-pointwise} measures how much the fitted mean prediction changes when the training set changes. It is related to epistemic or sampling uncertainty, but it is not a complete predictive-uncertainty measure.

For a future outcome, uncertainty includes at least:
\begin{equation}
  \underbrace{\operatorname{Var}_{\mathcal D}(\widehat f_{\mathcal D}(x))}_{\text{model or sampling variability}}
  +\underbrace{\sigma^2(x)}_{\text{outcome noise}},
\end{equation}
along with uncertainty caused by model misspecification and by estimating these terms. A stable but systematically wrong model can have low variance and poorly calibrated uncertainty.

If $B$ fitted predictors have common variance $v$ and pairwise correlation $\rho$, the variance of their average is
\begin{equation}
  \operatorname{Var}\left(\frac{1}{B}\sum_{b=1}^{B}\widehat f_b(x)\right)
  =v\left(\rho+\frac{1-\rho}{B}\right).
  \label{eq:ensemble-variance}
\end{equation}
To derive this identity, note that equal variances and pairwise correlation imply
\begin{equation*}
  \operatorname{Var}(\widehat f_b(x))=v,
  \qquad
  \operatorname{Cov}(\widehat f_b(x),\widehat f_{b'}(x))=\rho v
  \quad (b\neq b').
\end{equation*}
The variance of a sum contains both individual variances and covariances. There are $B$ variance terms and $B(B-1)/2$ unordered pairs, so
\begin{align*}
  \operatorname{Var}\left(\frac{1}{B}\sum_{b=1}^{B}\widehat f_b(x)\right)
  &=\frac{1}{B^2}
    \left[Bv+B(B-1)\rho v\right]\\
  &=v\left(\rho+\frac{1-\rho}{B}\right).
\end{align*}
The result can be read as
\begin{equation}
  \underbrace{v\rho}_{\text{shared component}}
  +
  \underbrace{\frac{v(1-\rho)}{B}}_{\text{unshared component reduced by averaging}}.
\end{equation}

If $\rho=0$, the models fluctuate independently and the ensemble variance is $v/B$. With $B=100$ and $v=1$, it falls to $0.01$. If $\rho=1$, all predictions move together and averaging changes nothing: the variance remains $v$. In the intermediate example $B=100$, $v=1$, and $\rho=0.2$, equation~\eqref{eq:ensemble-variance} gives
\begin{equation*}
  0.2+\frac{0.8}{100}=0.208.
\end{equation*}
The independent portion has nearly disappeared, but the common component creates a variance floor. Indeed, as $B\to\infty$, the variance approaches $v\rho$, not zero.

This explains bagging. Bootstrap samples differ, so unstable base learners such as deep decision trees can make partially different errors that average out. However, all samples originate from the same observed dataset and overlap substantially, so the fitted models remain correlated. Random forests introduce additional feature randomness partly to reduce this correlation. An ensemble benefits from individually useful predictors whose errors are not perfectly aligned; repeating an identical model many times provides no benefit.

Bootstrap refitting can also approximate the sampling distribution of predictions by repeatedly resampling the observed data and refitting the complete procedure. This is an approximation to hypothetical new samples from the population, not literal access to those samples. Its validity depends on the data structure, dependence assumptions, estimator, and sample size.

\section{Comparison of the sources}
\label{sec:sources}

\paragraph{Hastie, Tibshirani, and Friedman (2009).}
A reference textbook used here for the exact squared-loss decomposition, model assessment, the zero-one-loss warning, and cross-validation. Its main role is to establish the classical statistical framework. It is an authoritative scholarly synthesis, although that framework predates the modern double-descent literature.

\paragraph{Geman, Bienenstock, and Doursat (1992).}
A peer-reviewed journal article used for nonparametric estimation, sample-size dependence, smoothing, regularization, and representation. Its main contribution here is to frame bias and variance as a finite-information dilemma, rather than merely a choice between particular algorithms. It is foundational and mathematically careful, but its conclusions about neural networks are historically situated.

\paragraph{Belkin, Hsu, Ma, and Mandal (2019).}
A peer-reviewed research article used for the interpolation threshold, double descent, minimum-norm solutions, and implicit inductive bias. Its role is to extend the capacity axis beyond the classical optimum and distinguish raw parameter count from the complexity of the selected interpolant. It provides primary modern evidence and a unifying proposal, not a theorem that every learning system exhibits double descent.

\paragraph{scikit-learn documentation (accessed 2026).}
Maintained pedagogical and technical documentation used for concrete accounts of underfitting, overfitting, learning curves, validation curves, and practical levers. It is official and reproducible, but it is not a peer-reviewed theoretical source.

\paragraph{Domingos (2000).}
A peer-reviewed conference paper used for the loss-dependent generalization of bias and variance, including zero-one loss. It directly addresses the classification nuance, but its definitions differ from ordinary squared bias and variance.

\subsection{Two important differences among the sources}

\paragraph{Classical flexibility versus modern capacity.}
Geman et al. and Hastie et al. emphasize the familiar tendency for flexibility to increase variance. Belkin et al. do not deny the classical regime; they extend the capacity axis beyond interpolation and show that test risk can descend again. The difference concerns the relationship between the chosen capacity measure and the effective complexity of the selected solution.

\paragraph{Squared loss versus zero-one loss.}
Hastie et al. show that classification error is not the sum of squared bias and prediction variance. Domingos obtains a decomposition for zero-one loss by redefining the central prediction and the components in a loss-dependent way. These statements are compatible once the definitions are made explicit.

\section{Direct answers to the five synthesis questions}
\label{sec:answers}

\subsection*{1. What is the decomposition at a point $x$, and which term is irreducible?}

Under equation~\eqref{eq:data-model} and squared loss,
\begin{equation*}
  \mathbb E_{\mathcal D,\varepsilon}
  [(Y-\widehat f_{\mathcal D}(x))^2\mid X=x]
  =\left(f(x)-\ED[\widehat f_{\mathcal D}(x)]\right)^2
  +\VarD(\widehat f_{\mathcal D}(x))
  +\sigma^2(x).
\end{equation*}
The terms are squared systematic error, sensitivity to the training set, and outcome noise. The last term is irreducible for point prediction of $Y$.

\subsection*{2. What is random in the expectations?}

For bias and variance at a fixed $x$, redraw the training set of fixed size from the same population and refit the entire learning procedure. For total pointwise test error, also draw a fresh response noise independent of the training data. For integrated risk, additionally draw or average over the test input $X$. In a real application these repetitions remain the theoretical definition even though only one training dataset is observed. After conditioning on one already fitted model, the observed training set is fixed and only future test inputs and responses remain random.

\subsection*{3. Does the same decomposition exist for zero-one loss?}

Not with the ordinary squared-loss terms. Zero-one loss depends on threshold crossings, so squared bias and prediction variance interact rather than add to misclassification error. Loss-specific generalized decompositions exist, but their bias and variance definitions must be stated and should not be confused with the regression quantities.

\subsection*{4. What is double descent, and in which regime does it appear?}

Double descent is a test-risk curve that first follows the classical U shape, peaks near the interpolation threshold, and decreases again in the over-parameterized interpolation regime. To the right of the threshold, training error remains approximately zero while effective regularity can improve because the algorithm selects particular interpolating solutions.

\subsection*{5. Give two levers for variance and two for bias.}

To reduce variance: increase the amount of relevant training data; strengthen smoothing or regularization. To reduce bias: enlarge the useful hypothesis class; weaken excessive regularization or improve the representation. Each lever may worsen the other component, so it must be selected by independent validation.

\section{How the theory maps to the notebook exercises}
\label{sec:map}

\begin{description}
  \item[Exercise 1:] Treat the collection of fitted curves as an empirical distribution of $\widehat f_{\mathcal D}(x)$. Their center visualizes bias; their spread visualizes variance.
  \item[Exercise 2:] Implement equation~\eqref{eq:mcvar} and verify the averaged version of equation~\eqref{eq:bv-pointwise}. Use independent test noise and report a relative discrepancy.
  \item[Exercise 3:] Explain changes in the optimal degree through the stabilizing effect of larger $n$. Do not assume the degree must change monotonically in every finite simulation.
  \item[Exercise 4:] For $k$-NN, smaller $k$ means greater effective complexity. This direction is opposite to polynomial degree.
  \item[Exercise 5:] Cross-validation estimates total risk from one dataset. The distribution of selected degrees measures selection instability, not bias or variance directly.
  \item[Exercise 6:] Plot pointwise components. Expect high variance where the Beta design provides few observations, often near $x=1$. Replace constant $\sigma^2$ by $\sigma^2(0.5+2x)^2$.
  \item[Bonus:] Ridge should trace a bias-variance path as $\alpha$ changes. Very small $\alpha$ resembles unstable unregularized degree-15 regression; very large $\alpha$ overshrinks and raises bias.
\end{description}

\section{Common mistakes to avoid}

\begin{enumerate}
  \item Computing ``variance'' across test locations for one fitted model. The required variance is across refitted models at a fixed location.
  \item Comparing predictions with noisy $Y$ when estimating bias. Bias is measured relative to $f(x)$ in the simulation.
  \item Reusing training noise as test noise. Test error requires a fresh independent response.
  \item Adding constant $\sigma^2$ in the heteroscedastic exercise.
  \item Treating $k$ as increasing complexity in $k$-NN. Larger $k$ means more smoothing and usually lower complexity.
  \item Claiming that degree 15 alone demonstrates double descent. A second descent must actually be observed beyond interpolation.
  \item Equating low model variance with complete predictive certainty.
  \item Applying supervised preprocessing before cross-validation rather than separately inside each fold.
  \item Reporting only a selected hyperparameter without uncertainty, the validation curve, or its frequency across repeated datasets.
  \item Confusing an individual loss $L(Y,g(X))$ with its population expectation $\R(g)$ or with a finite test-set average $\widehat\R(g)$.
  \item Treating a fitted model from one observed dataset as if it were the same object as the learning procedure evaluated over repeated datasets.
  \item Assuming that real data remove the conceptual randomness of $\mathcal D$. They provide one realization of a sampling process; they do not make repeated-sampling bias and variance directly observable.
  \item Equating the number of parameters with effective complexity in every regime. Regularization and the algorithm's choice among many solutions also matter.
  \item Deleting polynomial terms solely because their fitted coefficients are numerically small, without accounting for scaling, collinearity, hierarchy, and validation.
  \item Interpreting hyperparameter instability as proof of predictive instability. Several degrees can be selected inconsistently while having nearly equal risk.
  \item Expecting infinitely many highly correlated ensemble members to eliminate variance. Their shared component remains as $B$ increases.
\end{enumerate}

\section{A balanced answer to the final discussion question}

The bias-variance framework remains a good guide because it separates systematic error, sensitivity to data, and irreducible noise, and because these ideas directly explain the assignment's simulations, regularization, nearest neighbors, ensembles, and the effect of sample size. It is especially powerful when the loss is squared error and the target function and data-generating process are controlled.

It is not a complete model-selection rule for modern machine learning. On real data, $f$ and $\sigma^2(x)$ are unknown; the relevant test distribution may shift; classification loss behaves differently; optimization and representation matter; and parameter count may poorly measure effective complexity. Double descent shows that interpolation is a regime to analyze, not an automatic diagnosis of harmful overfitting.

A defensible conclusion is therefore nuanced: use bias and variance as diagnostic concepts and as an exact decomposition when its assumptions apply, but choose models using validated predictive performance while accounting for regularization, optimization, data distribution, uncertainty, and the inductive bias of the full learning procedure.

\begin{thebibliography}{9}

\bibitem{hastie2009}
T. Hastie, R. Tibshirani, and J. Friedman.
\newblock \emph{The Elements of Statistical Learning: Data Mining, Inference, and Prediction}.
\newblock Springer Series in Statistics, second edition, 2009, Chapter 7, pp. 219--259.
\newblock \url{https://doi.org/10.1007/978-0-387-84858-7}.

\bibitem{geman1992}
S. Geman, E. Bienenstock, and R. Doursat.
\newblock Neural networks and the bias/variance dilemma.
\newblock \emph{Neural Computation}, 4(1):1--58, 1992.
\newblock \url{https://doi.org/10.1162/neco.1992.4.1.1}.

\bibitem{belkin2019}
M. Belkin, D. Hsu, S. Ma, and S. Mandal.
\newblock Reconciling modern machine-learning practice and the classical bias-variance trade-off.
\newblock \emph{Proceedings of the National Academy of Sciences}, 116(32):15849--15854, 2019.
\newblock \url{https://doi.org/10.1073/pnas.1903070116}.

\bibitem{sklearn2026}
The scikit-learn developers.
\newblock Validation curves: plotting scores to evaluate models.
\newblock scikit-learn user guide, accessed 30 September 2026.
\newblock \url{https://scikit-learn.org/stable/modules/learning_curve.html}.

\bibitem{domingos2000}
P. Domingos.
\newblock A unified bias-variance decomposition and its applications.
\newblock In \emph{Proceedings of the Seventeenth International Conference on Machine Learning}, pp. 231--238, 2000.
\newblock \url{https://homes.cs.washington.edu/~pedrod/bvd.pdf}.

\end{thebibliography}

\end{document}
